\chapter{Implications and Limitations}
\label{chap:discussion}

\section{Accuracy of the variance approximation}

The variance-sensitivity derivation describes how the estimated MI
variance changes between samples. Its accuracy and the accuracy of the
mean approximation contribute differently to the degrees of freedom.
At \(4\times4\), MI 0.02 and \(n=1{,}000\), the population
sampling-variance predictions are within about 3--4\% of the observed
values, while the mean estimates exceed the population values by about
17--19\%. Equation~\eqref{eq:redesign-df-decomposition} shows how squaring
the mean-ratio contribution accounts for most of the resulting difference
between the first-order and moment degrees of freedom. The derivation
can therefore describe the variability of the variance estimator
reasonably well before its two moment approximations jointly give
accurate degrees of freedom.

The distinction matters most near independence, where the population
first-order MI variance is small. The quadratic sampling contribution
then remains substantial relative to that population value. In
Appendix~\ref{sec:independence-boundary},
\(\Vhat\approx G/n\) and \(\widehat\nu\approx G/2\), with
\(G\) approximately following a noncentral chi-squared distribution.
At independence, the approximating chi-squared distribution has mean
\(d\). The local calculation therefore predicts mean plug-in degrees
of freedom near \(d/2\), while the variance-estimator moment target is
near \(d\). This explains why substituting empirical probabilities
into the first-order formula need not recover the sampling-moment target,
even when the local approximation describes the observed component means.

Away from independence, the scaled mean errors suggest a second-order
correction whose coefficient depends on the population table. A correction
near \(d/n\) describes much of the mean shift in the tested
\(4\times4\) populations, but the observed variation with MI and margins
requires a coefficient derived for the actual table. Such a correction
would address the numerator of the moment-matching expression. Its
estimated coefficient and its effect on the completed rejection rule
would still need evaluation.

\section{Calibration and power of the MI-specific rule}

Expanded Welch improves null rejection in the two selected small-sample
populations. For uniform \(4\times4\) tables at \(n=20\), it gives
0.04897, compared with 0.08352 for Wald and 0.07394 for the sample-size
comparator. For \(8\times8\) different-skew tables at \(n=100\), it
gives 0.05523, compared with 0.07466 and 0.07246. Every rule is valid in
all repetitions at these settings, so the reductions do not reflect
failed calculations. In the wider grid, however, Expanded Welch also
lowers already-conservative rates. The MI-specific cutoff improves
calibration in some populations without improving it uniformly.

The paired nominal-power comparison measures the cost of those cutoffs.
At two null SDs in the uniform \(4\times4\) setting, Expanded Welch
loses 2.986 percentage points relative to the sample-size comparator.
After matching null rejection empirically, it instead gains 3.629 points.
The two comparisons answer different questions: nominal power includes
the methods' different false-positive rates, while calibrated power
compares discrimination at approximately the same false-positive rate.
The latter gain indicates that the sample-dependent degrees of freedom
change the ordering of evidence usefully along this population path.
A fixed Student \(t\) reference cannot produce that gain after calibration,
because it has the same ordering by \(|T|\) as Wald.

The different-skew \(8\times8\) path also gives higher calibrated power
for Expanded Welch, but its rejection curve initially falls as the true
MI difference increases. At the null, the estimated difference has a
positive bias; increasing \(I(Q)\) makes the true signed difference
negative and initially moves the estimate towards zero. Empirical null
calibration changes the rejection threshold and leaves this movement
in the shared statistic. The calibrated-power advantage and the initial
decline are therefore both features of the tested comparison. At the
larger-sample \(4\times4\) setting, the nominal loss is small and the
calibrated differences lie well within the stated practical tolerance.

\section{Bias, uncertainty and use on observed data}

The sparse \(16\times16\) different-skew comparison demonstrates the
limit of changing the cutoff. At \(n=1{,}000\), all four rules reject
far above 5\%. At the 0.01-nat alternative, almost every rejection has
the wrong estimated ordering. The positive difference bias is larger
than the true negative difference, as shown in
Equation~\eqref{eq:redesign-alternative-bias-shift}. The four-cell
construction reproduces the failure, so it is not restricted to the
copying association pattern. The fixed moment reference also over-rejects:
fitting denominator moments accurately does not repair a shifted numerator.

The matched same-skew populations expose a separate uncertainty problem.
Their expected null difference bias cancels at equal sample sizes,
but the mean squared standard error is almost twice the observed
difference variance in the sparse copying setting. Rejection is strongly
conservative. With unequal sample sizes, even identical population tables
retain residual difference biases because each MI estimate has a different
sampling distribution. The allocation results consequently depend on
which population receives fewer observations, as well as on the
sample-dependent cutoff.

For an observed categorical comparison, report the MI difference, its
standard error and the p-value together. If Expanded Welch is used,
report its estimated degrees of freedom as well. The difference supplies
estimated size and direction, while the standard error and reference
distribution determine the uncertainty assigned to it. A small estimated
degree of freedom alone does not identify reliable calibration: it can
occur near independence where Wald already rejects too rarely.

Before choosing either rejection rule for an application, evaluate it on
plausible populations with the application's category definitions, sample
sizes and margin patterns. Include distinct equal-MI tables and relevant
alternatives, and assess false positives, power, estimated direction and
validity. This comparison should guide the choice before interpreting
the observed application's p-values. The online-shopping analysis
illustrates the interpretation of the estimates and cutoffs, while its
unknown population MI values preclude measuring calibration directly.

\section{Limits of the evidence}

The focus settings were selected using earlier calibration results.
The new samples estimate their performance independently, but selection
still limits the scope of the resulting power claims. Each focused curve
keeps \(I(P)=0.02\) and increases \(I(Q)\) along the copying family.
The study does not establish the same calibrated-power advantage for
reversed changes, different association locations or arbitrary equal-MI
population pairs. The broad grid supplies evidence about other sample
sizes, table sizes and margin regimes, rather than a universal guarantee.

The primary populations have identical row and column margins within
each table. The four-cell comparison changes where dependence is
concentrated while retaining MI and margins, but both constructions align
the association locations between populations. Broader evidence would
require unequal row and column margins, other association patterns and
additional allocations. Monte Carlo intervals concern the specified
populations; they do not measure performance variation over these
untested populations.

The moment reference and empirical calibration require known populations.
They are simulation benchmarks. Replacing Expanded Welch's sample-dependent
degrees of freedom by a fixed value changes both their magnitude and
their dependence on the statistic. The contrast does not isolate estimation
error. Likewise, correlations and distribution-shape summaries describe
the joint sampling behaviour without identifying separate causal effects.
The calibrated-power bootstrap includes fitting and evaluation uncertainty
for the fixed constructed alternatives, and its intervals are pointwise.

The mathematical setting assumes independent categorical observations,
fixed categories, positive population cell probabilities and positive
first-order MI variance. Paired samples require a covariance contribution
to the difference variance. Growing alphabets and continuous MI estimators
have different approximation errors. Exact independence requires the
second-order argument in Appendix~\ref{sec:independence-boundary}, and
the positive-MI sweep shows that proximity to that boundary also matters
at finite sample sizes.
